% ==============================================================================
% SECTION 04.04: CONTINUOUS UNIFORM DISTRIBUTION U(a, b)
% ==============================================================================

\subsection*{Problem Statements --- Section 04.04}

\subsubsection*{Fundamental Level}

\begin{problema}
	\label{prob:4.4.1}
	Consider the continuous uniform distribution $X \sim U(2, 8)$ with PDF $f(x) = \dfrac{1}{8-2} = \dfrac{1}{6}$ for $x \in [2, 8]$.
	\begin{enumerate}
		\item Verify that $\int_{2}^{8} f(x)\,dx = 1$.
		\item Compute the CDF $F(x)$ and qualitatively sketch its linear shape.
		\item Compute $P(3 \le X \le 6)$ and $P(X \ge 7)$.
	\end{enumerate}
\end{problema}

\begin{problema}
	\label{prob:4.4.2}
	Let $X \sim U(a, b)$ with parameters $a < b$.
	\begin{enumerate}
		\item Compute the expectation $\E[X] = \dfrac{a + b}{2}$ and the variance $\Var(X) = \dfrac{(b-a)^{2}}{12}$ using the integral definitions.
		\item Prove the linearity of expectation: if $Y = cX + d$, then $Y \sim U(ca + d, cb + d)$.
	\end{enumerate}
\end{problema}

\begin{problema}
	\label{prob:4.4.3}
	The waiting time for a bus at a stop (in minutes) is modeled as $X \sim U(0, 15)$ (uniform distribution between 0 and 15 minutes).
	\begin{enumerate}
		\item Compute the probability that the waiting time is less than 5 minutes.
		\item Compute the quantile $q_{0.75}$ (the time below which 75\% of passengers wait).
		\item If a passenger has waited 10 minutes without seeing the bus, what is the conditional probability that they must wait at least 5 more minutes? Use the memorylessness property.
	\end{enumerate}
\end{problema}

\subsubsection*{Operational Level}

\begin{problema}
	\label{prob:4.4.4}
	Let $X \sim U(0, 1)$. Explicitly compute the quantiles $q_{0.10}$, $q_{0.50}$ (median), and $q_{0.90}$.
\end{problema}

\begin{problema}
	\label{prob:4.4.5}
	Implement a Monte Carlo simulation of 100{,}000 samples from $U(0, 1)$ and verify that the empirical mean and variance match $\E[X] = 0.5$ and $\Var(X) = 1/12 \approx 0.0833$. Additionally, plot the histogram with the PDF overlaid.
\end{problema}

\begin{problema}
	\label{prob:4.4.6}
	In Monte Carlo simulation, we want to generate samples of a variable $Y$ with an arbitrary CDF $F(y)$. Prove that if $U \sim U(0, 1)$ and $F$ is invertible, then $Y = F^{-1}(U)$ has CDF $F$. Apply this result to generate 50{,}000 samples of $Y \sim \text{Exp}(\lambda = 2)$ from the uniform distribution.
\end{problema}

\subsubsection*{Analytical Level}

\begin{problema}
	\label{prob:4.4.7}
	Let $X$ be a continuous variable with bounded support $[a, b]$. Prove that the uniform distribution $U(a, b)$ \emph{maximizes the differential entropy} $h(X) = -\int f(x) \log f(x)\,dx$ among all distributions with support in $[a, b]$. Hint: use Lagrange multipliers with the constraint $\int f(x)\,dx = 1$.
\end{problema}

\begin{problema}
	\label{prob:4.4.8}
	Let $X_{1}, X_{2}, \dots, X_{n}$ be i.i.d. samples from $U(0, 1)$ and $X_{(1)} \le X_{(2)} \le \dots \le X_{(n)}$ their order statistics.
	\begin{enumerate}
		\item Prove that $X_{(k)}$ has PDF Beta($k, n - k + 1$): $f_{X_{(k)}}(x) = \dfrac{n!}{(k-1)!(n-k)!} x^{k-1} (1-x)^{n-k}$ for $x \in [0, 1]$.
		\item Compute $\E[X_{(k)}] = \dfrac{k}{n+1}$ and $\Var(X_{(k)})$.
	\end{enumerate}
\end{problema}

\subsubsection*{Challenging Level}

\begin{problema}
	\label{prob:4.4.9}
	Let $X_{1}, \dots, X_{n}$ be i.i.d. from $U(\theta - 1/2, \theta + 1/2)$ where $\theta$ is unknown. The support has length 1. Use the maximum likelihood method to estimate $\theta$. Is the estimator $\hat{\theta}$ unbiased? What is its asymptotic variance?
\end{problema}

\begin{problema}
	\label{prob:4.4.10}
	The Kolmogorov-Smirnov test compares the empirical CDF $F_{n}$ with the theoretical CDF $F_{0}$. For $n = 50$ observations suspected to come from $U(0, 1)$ with $D_{50} = 0.18$:
	\begin{enumerate}
		\item Compute the approximate $p$-value: $p \approx 2 \sum_{k=1}^{\infty}(-1)^{k-1} e^{-2k^{2}n D_{n}^{2}}$.
		\item Compare with the critical value $K_{0.95}/\sqrt{n} = 1.358/\sqrt{50} \approx 0.192$.
		\item Conclude whether the data are consistent with $U(0, 1)$ at level $\alpha = 0.05$.
	\end{enumerate}
\end{problema}

\subsection*{Hints for the Problems --- Section 04.04}

\begin{sugerencia}
	\label{sug:4.4.1}
	For Problem \ref{prob:4.4.1}, $\int_{2}^{8} (1/6)\,dx = 6/6 = 1$. The CDF is $F(x) = (x-2)/6$ for $x \in [2, 8]$. $P(3 \le X \le 6) = (6-2)/6 = 0.5$. $P(X \ge 7) = (8-7)/6 = 1/6 \approx 0.167$.
\end{sugerencia}

\begin{sugerencia}
	\label{sug:4.4.2}
	For Problem \ref{prob:4.4.2}, $\E[X] = \int_{a}^{b} x \cdot (1/(b-a))\,dx = (b^{2} - a^{2})/(2(b-a)) = (a+b)/2$. $\E[X^{2}] = (b^{2} + ab + a^{2})/3$. $\Var(X) = (b-a)^{2}/12$. For $Y = cX + d$: the PDF of $Y$ is $f_{Y}(y) = f_{X}((y-d)/c) \cdot (1/|c|) = 1/((cb+d)-(ca+d)) = 1/(c(b-a))$ on $[ca+d, cb+d]$.
\end{sugerencia}

\begin{sugerencia}
	\label{sug:4.4.3}
	For Problem \ref{prob:4.4.3}, $F(x) = x/15$ for $x \in [0, 15]$. $P(X < 5) = 5/15 = 1/3 \approx 0.333$. $q_{0.75}$: $x/15 = 0.75$, giving $x = 11.25$ min. Memorylessness property: $P(X \ge 15 \mid X \ge 10) = P(X \ge 5) = 10/15 = 2/3 \approx 0.667$. By symmetry, $P(X \ge a + b \mid X \ge a) = P(X \ge b)$.
\end{sugerencia}

\begin{sugerencia}
	\label{sug:4.4.4}
	For Problem \ref{prob:4.4.4}, $q_{p}$ for $U(0, 1)$ is $q_{p} = p$. Thus $q_{0.10} = 0.1$, $q_{0.50} = 0.5$ (median), $q_{0.90} = 0.9$.
\end{sugerencia}

\begin{sugerencia}
	\label{sug:4.4.5}
	For Problem \ref{prob:4.4.5}, use `np.random.seed(42)` and `np.random.uniform(0, 1, 100000)`. Verify with `np.mean(samples)` ($\approx 0.5$) and `np.var(samples, ddof=1)` ($\approx 0.0833$). For the histogram, use `plt.hist(samples, bins=50, density=True)` and overlay the constant PDF $1$.
\end{sugerencia}

\begin{sugerencia}
	\label{sug:4.4.6}
	For Problem \ref{prob:4.4.6}, $P(Y \le y) = P(F^{-1}(U) \le y) = P(U \le F(y)) = F(y)$ by monotonicity of $F^{-1}$ and uniformity of $U$. For Exponential($\lambda = 2$): $F(y) = 1 - e^{-2y}$ for $y \ge 0$, so $Y = -\ln(1-U)/2$.
\end{sugerencia}

\begin{sugerencia}
	\label{sug:4.4.7}
	For Problem \ref{prob:4.4.8}, maximize $h[X] = -\int f(x) \log f(x)\,dx$ subject to $\int f(x)\,dx = 1$ and support in $[a, b]$. Lagrangian: $\mathcal{L} = -\int f \log f\,dx - \lambda(\int f\,dx - 1)$. Differentiating with respect to $f$: $-\log f - 1 - \lambda = 0$, so $f(x) = e^{-1-\lambda} = $ constant. The constant is $1/(b-a)$ by normalization. Therefore $U(a,b)$ maximizes the entropy.
\end{sugerencia}

\begin{sugerencia}
	\label{sug:4.4.8}
	For Problem \ref{prob:4.4.8}, the PDF of $X_{(k)}$ is Beta($k, n-k+1$): $f_{X_{(k)}}(x) = \frac{n!}{(k-1)!(n-k)!} x^{k-1}(1-x)^{n-k}$. The expectation is $\E[X_{(k)}] = k/(n+1)$ (Beta symmetry). The variance is $\Var(X_{(k)}) = \frac{k(n-k+1)}{(n+1)^{2}(n+2)}$.
\end{sugerencia}

\begin{sugerencia}
	\label{sug:4.4.9}
	For Problem \ref{prob:4.4.9}, the log-likelihood is $\ell(\theta) = \sum \log f(x_i \mid \theta) = -n \log(1) = 0$ for any $\theta$ in the admissible range. The likelihood is constant: any $\hat{\theta}$ that covers all observations is an MLE. The natural estimator is $\hat{\theta} = \frac{1}{2}(\max X_i + \min X_i)$ (midrange). Bias: $\E[\hat{\theta}] = \theta$ (unbiased). Variance: $\Var(\hat{\theta}) = 1/(2(n+1)(n+2))$ (from the variance of the range).
\end{sugerencia}

\begin{sugerencia}
	\label{sug:4.4.10}
	For Problem \ref{prob:4.4.10}, with $n = 50$ and $D_{50} = 0.18$, $nD_{n}^{2} = 50 \cdot 0.0324 = 1.62$. The $p$-value is $p \approx 2 \sum_{k=1}^{\infty}(-1)^{k-1} e^{-2k^{2} \cdot 1.62}$. For $k=1$: $2 e^{-3.24} \approx 2 \cdot 0.0392 = 0.0784$. For $k=2$: $-2 e^{-12.96} \approx 0$. Therefore $p \approx 0.078$. Since $D_{50} = 0.18 < 0.192 = 1.358/\sqrt{50}$, we do not reject $H_{0}$ at level $\alpha = 0.05$.
\end{sugerencia}

\subsection*{Solutions to the Problems --- Section 04.04}

\begin{solucion}
	\label{sol:4.4.1}
	For Problem \ref{prob:4.4.1}:
	\begin{enumerate}
		\item $\int_{2}^{8} \frac{1}{6}\,dx = \frac{1}{6}(8-2) = 1$. Verification completed.
		\item The CDF is linear:
		\begin{align*}
			F(x) = P(X \le x) = \int_{2}^{x} \frac{1}{6}\,dt = \frac{x-2}{6}, \quad x \in [2, 8].
		\end{align*}
		\item $P(3 \le X \le 6) = F(6) - F(3) = \frac{4}{6} - \frac{1}{6} = \frac{3}{6} = 0.5$. $P(X \ge 7) = 1 - F(7) = 1 - \frac{5}{6} = \frac{1}{6} \approx 0.167$.
	\end{enumerate}
\end{solucion}

\begin{solucion}
	\label{sol:4.4.2}
	For Problem \ref{prob:4.4.2}:
	\begin{enumerate}
		\item $\E[X] = \int_{a}^{b} x \cdot \frac{1}{b-a}\,dx = \frac{1}{b-a} \cdot \frac{x^{2}}{2}\Big|_{a}^{b} = \frac{b^{2} - a^{2}}{2(b-a)} = \frac{a+b}{2}$. $\E[X^{2}] = \int_{a}^{b} x^{2} \cdot \frac{1}{b-a}\,dx = \frac{b^{3} - a^{3}}{3(b-a)} = \frac{a^{2} + ab + b^{2}}{3}$. $\Var(X) = \E[X^{2}] - (\E[X])^{2} = \frac{a^{2} + ab + b^{2}}{3} - \frac{(a+b)^{2}}{4} = \frac{(b-a)^{2}}{12}$.
		\item For $Y = cX + d$ with $c > 0$: using the change of variable $x = (y-d)/c$, $dx = dy/c$, and support $[a, b]$ mapped to $[ca+d, cb+d]$:
		\begin{align*}
			f_{Y}(y) = f_{X}\left(\frac{y-d}{c}\right) \cdot \frac{1}{c} = \frac{1}{b-a} \cdot \frac{1}{c} = \frac{1}{(cb+d) - (ca+d)} = \frac{1}{c(b-a)},
		\end{align*}
		confirming $Y \sim U(ca+d, cb+d)$.
	\end{enumerate}
\end{solucion}

\begin{solucion}
	\label{sol:4.4.3}
	For Problem \ref{prob:4.4.3} with $X \sim U(0, 15)$:
	\begin{enumerate}
		\item $P(X < 5) = F(5) = 5/15 = 1/3 \approx 0.333$.
		\item $q_{0.75}$ satisfies $F(q_{0.75}) = q_{0.75}/15 = 0.75$, giving $q_{0.75} = 11.25$ minutes. Therefore, 75\% of passengers wait less than 11.25 minutes.
		\item The uniform distribution does NOT have the memorylessness property (unlike the exponential). The conditional probability is $P(X \ge 15 \mid X \ge 10) = \frac{P(X \ge 15)}{P(X \ge 10)} = \frac{0/15}{5/15} = 0$. That is, if 10 minutes have passed without a bus, it is because there is NO bus in the first 10 minutes, so the probability of waiting at least 5 more minutes is 0. \emph{Correction:} the uniform does NOT have memorylessness, so $P(X \ge 15 \mid X \ge 10) = P(X \ge 5) = 10/15 = 2/3 \approx 0.667$ (this would be correct only if the uniform \emph{did} have memorylessness, which is NOT the case). The correct calculation: conditionally on $X \ge 10$, $X$ follows $U(10, 15)$ (by the \emph{random inspection} property), so $P(X \ge 15 \mid X \ge 10) = 0$ (since the support only extends up to 15).
	\end{enumerate}
\end{solucion}

\begin{solucion}
	\label{sol:4.4.4}
	For Problem \ref{prob:4.4.4} with $X \sim U(0, 1)$: the CDF is $F(x) = x$ for $x \in [0, 1]$. Therefore, $q_{p} = F^{-1}(p) = p$. In particular:
	\begin{itemize}
		\item $q_{0.10} = 0.1$ (10th percentile)
		\item $q_{0.50} = 0.5$ (median)
		\item $q_{0.90} = 0.9$ (90th percentile)
	\end{itemize}
	The quantiles of $U(0, 1)$ are trivially equal to the probabilities.
\end{solucion}

\begin{solucion}
	\label{sol:4.4.5}
	For Problem \ref{prob:4.4.5} with `np.random.seed(42)` and `np.random.uniform(0, 1, 100000)`:
	\begin{align*}
		\bar{X} \approx 0.5000 \quad (\text{theoretical: } 0.5), \quad s^{2} \approx 0.0833 \quad (\text{theoretical: } 1/12 \approx 0.08333).
	\end{align*}
	The histogram with 50 bins shows a uniform distribution with density $\approx 1$ on $[0, 1]$, confirming the theory.
\end{solucion}

\begin{solucion}
	\label{sol:4.4.6}
	For Problem \ref{prob:4.4.6}, we prove the inversion method:
	\begin{align*}
		P(Y \le y) = P(F^{-1}(U) \le y) = P(U \le F(y)) = F(y)
	\end{align*}
	by the monotonicity of $F^{-1}$ and the uniformity of $U$. Therefore, $Y$ has CDF $F$. For Exponential($\lambda = 2$): $F(y) = 1 - e^{-2y}$ for $y \ge 0$, so $Y = -\frac{1}{2}\ln(1-U)$. With `np.random.seed(42)` and `np.random.uniform(0, 1, 50000)`, the 50{,}000 samples of $Y$ have empirical mean $\approx 0.5$ and variance $\approx 0.25$, consistent with Exponential(2).
\end{solucion}

\begin{solucion}
	\label{sol:4.4.7}
	For Problem \ref{prob:4.4.7}, we maximize the differential entropy $h[X] = -\int_{a}^{b} f(x) \log f(x)\,dx$ subject to $\int_{a}^{b} f(x)\,dx = 1$. The Lagrangian is:
	\begin{align*}
		\mathcal{L}[f] = -\int_{a}^{b} f(x) \log f(x)\,dx - \lambda \left( \int_{a}^{b} f(x)\,dx - 1 \right).
	\end{align*}
	Using calculus of variations (functional derivative with respect to $f$): $-\log f(x) - 1 - \lambda = 0$, which gives $\log f(x) = -(1 + \lambda)$, that is, $f(x) = e^{-(1+\lambda)}$ is constant. Normalization requires $f(x) = 1/(b-a)$ on $[a, b]$. Therefore, $U(a, b)$ maximizes the entropy among all distributions with support in $[a, b]$. This justifies why the uniform is the "least informative" prior in Bayesian inference. \quad \qedhere
\end{solucion}

\begin{solucion}
	\label{sol:4.4.8}
	For Problem \ref{prob:4.4.8}:
	\begin{enumerate}
		\item The PDF of $X_{(k)}$ (the $k$-th order statistic of $n$ samples from $U(0,1)$) is obtained by noting that exactly $k-1$ samples are $< x$, exactly one falls in $[x, x+dx]$, and exactly $n-k$ are $> x+dx$. The number of ways to choose these samples is $\binom{n}{k-1, 1, n-k} = \frac{n!}{(k-1)!(n-k)!}$. Thus:
		\begin{align*}
			f_{X_{(k)}}(x) = \frac{n!}{(k-1)!(n-k)!} x^{k-1}(1-x)^{n-k} \cdot \mathbb{1}_{[0,1]}(x),
		\end{align*}
		which is exactly the PDF of a Beta($k, n-k+1$) distribution.
		\item For Beta($\alpha, \beta$), $\E[X] = \alpha/(\alpha+\beta)$. With $\alpha = k$, $\beta = n-k+1$: $\E[X_{(k)}] = k/(n+1)$. The variance is $\Var(X_{(k)}) = \frac{\alpha\beta}{(\alpha+\beta)^{2}(\alpha+\beta+1)} = \frac{k(n-k+1)}{(n+1)^{2}(n+2)}$. For $k=1$ (minimum) and $k=n$ (maximum), the variance reduces to $n/((n+1)^{2}(n+2))$, reflecting the concentration at the extremes.
	\end{enumerate}
\end{solucion}

\begin{solucion}
	\label{sol:4.4.9}
	For Problem \ref{prob:4.4.9} with $X_{i} \sim U(\theta - 1/2, \theta + 1/2)$:
	\begin{enumerate}
		\item The log-likelihood is $\ell(\theta) = \sum \log f(x_i \mid \theta)$. Since $f(x \mid \theta) = 1$ for $x \in [\theta - 1/2, \theta + 1/2]$ and $0$ otherwise, the log-likelihood is $0$ if all observations lie within the interval $[\theta - 1/2, \theta + 1/2]$, and $-\infty$ otherwise. The MLE is any $\hat{\theta}$ such that $\max X_i \le \hat{\theta} + 1/2$ and $\min X_i \ge \hat{\theta} - 1/2$, that is, $\hat{\theta} \in [\max X_i - 1/2, \min X_i + 1/2]$. The natural estimator is the \emph{midrange}: $\hat{\theta} = (X_{(1)} + X_{(n)})/2$.
		\item Bias: $\E[\hat{\theta}] = (\E[X_{(1)}] + \E[X_{(n)}])/2 = (1/(n+1) + n/(n+1))/2 = 1/2 = \theta$, so $\hat{\theta}$ is unbiased. Variance: $\Var(\hat{\theta}) = (\Var(X_{(1)}) + \Var(X_{(n)}) + 2\cov(X_{(1)}, X_{(n)}))/4$. For large $n$, $\Var(\hat{\theta}) \approx 1/(2(n+1)(n+2)) \to 0$ at rate $1/n^{2}$, showing that the midrange is highly efficient for the uniform.
	\end{enumerate}
\end{solucion}

\begin{solucion}
	\label{sol:4.4.10}
	For Problem \ref{prob:4.4.10} with $n = 50$ and $D_{50} = 0.18$:
	\begin{enumerate}
		\item The $p$-value of the KS test is approximated by the series:
		\begin{align*}
			p \approx 2 \sum_{k=1}^{\infty}(-1)^{k-1} e^{-2k^{2} n D_{n}^{2}} = 2 \sum_{k=1}^{\infty}(-1)^{k-1} e^{-2k^{2} \cdot 1.62}.
		\end{align*}
		For $k=1$: $2 e^{-3.24} \approx 0.0784$. For $k=2$: $-2 e^{-12.96} \approx -2.3 \times 10^{-6} \approx 0$. Therefore, $p \approx 0.078$.
		\item Critical value: $K_{0.95}/\sqrt{n} = 1.358/\sqrt{50} \approx 0.192$. Since $D_{50} = 0.18 < 0.192$, we do not reject $H_{0}$ at level $\alpha = 0.05$.
		\item Conclusion: with $p \approx 0.078 > 0.05$ and $D_{50} < K_{0.95}/\sqrt{n}$, there is no significant statistical evidence to reject that the data come from $U(0, 1)$. The KS test does not detect significant deviations at the 5\% level. \quad \qedhere
	\end{enumerate}
\end{solucion}

% ==============================================================================
% SECTION 04.05: EXPONENTIAL DISTRIBUTION AND MEMORYLESS PROCESSES
% ==============================================================================

\subsection*{Problem Statements --- Section 04.05}

\subsubsection*{Fundamental Level}

\begin{problema}
	\label{prob:4.5.1}
	A continuous random variable $X$ has an exponential distribution with parameter $\lambda = 2$ and PDF $f(x) = 2e^{-2x}$ for $x \ge 0$.
	\begin{enumerate}
		\item Verify that $\int_{0}^{\infty} f(x)\,dx = 1$.
		\item Compute the CDF $F(x) = 1 - e^{-2x}$ for $x \ge 0$.
		\item Compute $P(X \le 1)$ and $P(X \ge 3)$.
	\end{enumerate}
\end{problema}

\begin{problema}
	\label{prob:4.5.2}
	Let $X \sim \text{Exp}(\lambda)$ with $\lambda > 0$.
	\begin{enumerate}
		\item Compute the expectation $\E[X] = 1/\lambda$ and the variance $\Var(X) = 1/\lambda^{2}$ by integration.
		\item Verify that the \emph{moment generating function} is $M_{X}(t) = \lambda/(\lambda - t)$ for $t < \lambda$.
	\end{enumerate}
\end{problema}

\begin{problema}
	\label{prob:4.5.3}
	In a Poisson process with rate $\lambda = 3$ events per hour, let $X$ be the time until the first event.
	\begin{enumerate}
		\item Prove that $X \sim \text{Exp}(3)$ with PDF $f(x) = 3e^{-3x}$ for $x \ge 0$.
		\item Compute the mean time between events and its standard deviation.
		\item Compute the probability that the first event occurs after 30 minutes.
	\end{enumerate}
\end{problema}

\subsubsection*{Operational Level}

\begin{problema}
	\label{prob:4.5.4}
	Formally prove the \emph{memorylessness property} of the exponential distribution: if $X \sim \text{Exp}(\lambda)$, then for all $s, t > 0$:
	$$P(X > s + t \mid X > s) = P(X > t)$$
	\begin{enumerate}
		\item Use the definition of conditional probability and the form of the exponential CDF.
		\item Show that the exponential is the \emph{only} continuous distribution with this property (among distributions with support in $[0, \infty)$).
	\end{enumerate}
\end{problema}

\begin{problema}
	\label{prob:4.5.5}
	Let $X \sim \text{Exp}(\lambda = 0.5)$ (mean 2 minutes).
	\begin{enumerate}
		\item Compute the quantiles $q_{0.10}$, $q_{0.50}$ (median), and $q_{0.90}$.
		\item Find the quantile $q_{0.95}$ and explain its meaning in the context of a queueing system.
	\end{enumerate}
\end{problema}

\begin{problema}
	\label{prob:4.5.6}
	Implement a Monte Carlo simulation of 100{,}000 samples from $\text{Exp}(\lambda = 2)$ and verify that the empirical mean and variance match $\E[X] = 0.5$ and $\Var(X) = 0.25$. Additionally, plot the histogram with the PDF overlaid.
\end{problema}

\subsubsection*{Analytical Level}

\begin{problema}
	\label{prob:4.5.7}
	Let $X_{1}, \dots, X_{n}$ be i.i.d. from $\text{Exp}(\lambda)$. Derive the maximum likelihood estimator (MLE) of the parameter $\lambda$.
	\begin{enumerate}
		\item Write the log-likelihood and differentiate with respect to $\lambda$.
		\item Prove that $\hat{\lambda} = 1/\bar{X}$ where $\bar{X}$ is the sample mean.
		\item Compute the bias and asymptotic variance of $\hat{\lambda}$.
	\end{enumerate}
\end{problema}

\begin{problema}
	\label{prob:4.5.8}
	Let $X_{1}, \dots, X_{n}$ be i.i.d. from $\text{Exp}(\lambda)$.
	\begin{enumerate}
		\item Prove that the sum $S_{n} = \sum X_{i}$ follows the \emph{Erlang} distribution (or Gamma with integer shape parameter): $S_{n} \sim \text{Gamma}(n, \lambda)$.
		\item Compute $\E[S_{n}]$ and $\Var(S_{n})$.
		\item For $n = 5$ and $\lambda = 1$, plot the PDF of $S_{5}$ and compare with the Normal approximation via the CLT.
	\end{enumerate}
\end{problema}

\subsubsection*{Challenging Level}

\begin{problema}
	\label{prob:4.5.9}
	In a reliability system with $n$ components in series, the system fails when the first component fails. If the times to failure are i.i.d. $X_{i} \sim \text{Exp}(\lambda)$ independent, the time to system failure is $T = \min(X_{1}, \dots, X_{n})$.
	\begin{enumerate}
		\item Prove that $T \sim \text{Exp}(n\lambda)$.
		\item Compute the mean time to system failure.
		\item For $n = 10$ components with $\lambda = 0.001$ failures/hour, compute $P(T > 100)$ and the median failure time.
	\end{enumerate}
\end{problema}

\begin{problema}
	\label{prob:4.5.10}
	In a service station with Poisson($\lambda$) arrivals and $\text{Exp}(\mu)$ service times:
	\begin{enumerate}
		\item Prove that the interarrival time is $\text{Exp}(\lambda)$ and that a customer's total service time is $\text{Exp}(\mu)$.
		\item Compute the traffic intensity $\rho = \lambda/\mu$ and explain its meaning (the system is stable if $\rho < 1$).
		\item For $\lambda = 5$ customers/min and $\mu = 6$ customers/min, compute the probability that the system is busy at a random instant.
	\end{enumerate}
\end{problema}

\subsection*{Hints for the Problems --- Section 04.05}

\begin{sugerencia}
	\label{sug:4.5.1}
	For Problem \ref{prob:4.5.1}, $\int_{0}^{\infty} 2e^{-2x}\,dx = 2 \cdot (1/2) = 1$. The CDF is $F(x) = 1 - e^{-2x}$. $P(X \le 1) = 1 - e^{-2} \approx 0.865$. $P(X \ge 3) = e^{-6} \approx 0.00248$.
\end{sugerencia}

\begin{sugerencia}
	\label{sug:4.5.2}
	For Problem \ref{prob:4.5.2}, $\E[X] = \int_{0}^{\infty} x \lambda e^{-\lambda x}\,dx = 1/\lambda$ (by integration by parts). $\E[X^{2}] = 2/\lambda^{2}$. $\Var(X) = 2/\lambda^{2} - 1/\lambda^{2} = 1/\lambda^{2}$. The MGF is $M_{X}(t) = \int_{0}^{\infty} e^{tx} \lambda e^{-\lambda x}\,dx = \lambda/(\lambda - t)$ for $t < \lambda$.
\end{sugerencia}

\begin{sugerencia}
	\label{sug:4.5.3}
	For Problem \ref{prob:4.5.3}, by the property of the Poisson process, $P(X > x) = e^{-3x}$ (probability of zero events in time $x$). Thus $F(x) = 1 - e^{-3x}$ and $f(x) = 3e^{-3x}$. $\E[X] = 1/3$ hour = 20 min. $\sigma = 1/3 \approx 0.333$ h. $P(X \ge 0.5) = e^{-1.5} \approx 0.223$.
\end{sugerencia}

\begin{sugerencia}
	\label{sug:4.5.4}
	For Problem \ref{prob:4.5.4}, $P(X > s + t \mid X > s) = P(X > s + t)/P(X > s) = e^{-\lambda(s+t)}/e^{-\lambda s} = e^{-\lambda t} = P(X > t)$. For uniqueness: if $g$ is the tail $P(X > x)$, the property implies $g(s+t) = g(s) g(t)$. The only measurable solutions with $g(0) = 1$ are $g(x) = e^{-cx}$ for some constant $c > 0$. Thus the exponential is unique.
\end{sugerencia}

\begin{sugerencia}
	\label{sug:4.5.5}
	For Problem \ref{prob:4.5.5}, $q_{p} = -\frac{1}{0.5}\ln(1-p) = -2\ln(1-p)$. $q_{0.10} \approx 0.211$, $q_{0.50} = 2 \ln 2 \approx 1.386$, $q_{0.90} \approx 4.605$ min. The quantile $q_{0.95} \approx 5.991$ min means that 95\% of customers wait less than 6 min.
\end{sugerencia}

\begin{sugerencia}
	\label{sug:4.5.6}
	For Problem \ref{prob:4.5.6}, use `np.random.seed(42)` and `np.random.exponential(1/2, 100000)`. Verify with `np.mean(samples)` ($\approx 0.5$) and `np.var(samples, ddof=1)` ($\approx 0.25$). The histogram shows a distribution with a heavy right tail, decaying exponentially.
\end{sugerencia}

\begin{sugerencia}
	\label{sug:4.5.7}
	For Problem \ref{prob:4.5.7}, the log-likelihood is $\ell(\lambda) = n\log \lambda - \lambda \sum x_i$. Differentiating: $\ell'(\lambda) = n/\lambda - \sum x_i = 0$, giving $\hat{\lambda} = n/\sum x_i = 1/\bar{x}$. Bias: $\E[\hat{\lambda}] = n\lambda/(n-1) = \lambda \cdot n/(n-1)$, so $\hat{\lambda}$ is biased. The asymptotic variance is $\lambda^{2}/n$ (via Fisher information $I(\lambda) = 1/\lambda^{2}$).
\end{sugerencia}

\begin{sugerencia}
	\label{sug:4.5.8}
	For Problem \ref{prob:4.5.8}, the PDF of $S_{n}$ is obtained by repeated convolution or via the MGF: $M_{S_{n}}(t) = (\lambda/(\lambda - t))^{n}$ for $t < \lambda$, which is the MGF of $\text{Gamma}(n, \lambda)$. By the uniqueness of the MGF, $S_{n} \sim \text{Gamma}(n, \lambda)$ with $\E[S_{n}] = n/\lambda$ and $\Var(S_{n}) = n/\lambda^{2}$. For $n = 5$, $\lambda = 1$: $\E[S_{5}] = 5$, $\Var(S_{5}) = 5$.
\end{sugerencia}

\begin{sugerencia}
	\label{sug:4.5.9}
	For Problem \ref{prob:4.5.9}, $P(T > t) = P(\min X_{i} > t) = \prod P(X_{i} > t) = e^{-n\lambda t}$ by independence. Thus $T \sim \text{Exp}(n\lambda)$. $\E[T] = 1/(n\lambda)$. With $n = 10$ and $\lambda = 0.001$: $\E[T] = 100$ h, $P(T > 100) = e^{-1} \approx 0.368$, median $= 100 \ln 2 \approx 69.3$ h.
\end{sugerencia}

\begin{sugerencia}
	\label{sug:4.5.10}
	For Problem \ref{prob:4.5.10}, in an $M/M/1$ queueing system, the interarrival time is $\text{Exp}(\lambda)$ and the service time is $\text{Exp}(\mu)$ by the properties of the Poisson process. The traffic intensity is $\rho = \lambda/\mu$ (fraction of time the server is busy). By the stationary distribution, the probability that the system is busy is $\rho = \lambda/\mu$. With $\lambda = 5$ and $\mu = 6$: $\rho = 5/6 \approx 0.833$, so the server is busy 83.3\% of the time.
\end{sugerencia}

\subsection*{Solutions to the Problems --- Section 04.05}

\begin{solucion}
	\label{sol:4.5.1}
	For Problem \ref{prob:4.5.1}:
	\begin{enumerate}
		\item $\int_{0}^{\infty} 2e^{-2x}\,dx = 2 \cdot [-e^{-2x}/2]_{0}^{\infty} = 2 \cdot (0 - (-1/2)) = 1$. Verification completed.
		\item $F(x) = \int_{0}^{x} 2e^{-2t}\,dt = 1 - e^{-2x}$ for $x \ge 0$.
		\item $P(X \le 1) = 1 - e^{-2} \approx 1 - 0.1353 = 0.8647$. $P(X \ge 3) = e^{-6} \approx 0.00248$.
	\end{enumerate}
\end{solucion}

\begin{solucion}
	\label{sol:4.5.2}
	For Problem \ref{prob:4.5.2}:
	\begin{enumerate}
		\item By integration by parts with $u = x$, $dv = \lambda e^{-\lambda x}\,dx$:
		\begin{align*}
			\E[X] &= \int_{0}^{\infty} x \lambda e^{-\lambda x}\,dx = \left[ -x e^{-\lambda x} \right]_{0}^{\infty} + \int_{0}^{\infty} e^{-\lambda x}\,dx = 0 + \frac{1}{\lambda} = \frac{1}{\lambda}.
		\end{align*}
		$\E[X^{2}] = 2/\lambda^{2}$ (by induction or double integration by parts). $\Var(X) = 2/\lambda^{2} - 1/\lambda^{2} = 1/\lambda^{2}$. \emph{Correction:} $\E[X^{2}]$ for the Exponential is $2/\lambda^{2}$ (not $2 \cdot 1/\lambda^{2}$), so $\Var(X) = 2/\lambda^{2} - 1/\lambda^{2} = 1/\lambda^{2}$. The standard deviation is $\sigma = 1/\lambda$.
		\item $M_{X}(t) = \int_{0}^{\infty} e^{tx} \lambda e^{-\lambda x}\,dx = \lambda \int_{0}^{\infty} e^{-(\lambda - t)x}\,dx = \dfrac{\lambda}{\lambda - t}$ for $t < \lambda$.
	\end{enumerate}
\end{solucion}

\begin{solucion}
	\label{sol:4.5.3}
	For Problem \ref{prob:4.5.3} (Poisson process with $\lambda = 3$):
	\begin{enumerate}
		\item In a Poisson process with rate $\lambda$, the number of events in time $t$ is $N(t) \sim \text{Poisson}(\lambda t)$. The time until the first event satisfies $P(X > x) = P(N(x) = 0) = e^{-\lambda x}$. Differentiating: $f(x) = \lambda e^{-\lambda x}$, that is, $X \sim \text{Exp}(\lambda = 3)$.
		\item $\E[X] = 1/3$ hour = 20 minutes. $\sigma = 1/3$ hour = 20 minutes. The mean and standard deviation are equal in the exponential.
		\item $P(X \ge 0.5) = e^{-3 \cdot 0.5} = e^{-1.5} \approx 0.223$. There is a 22.3\% probability that the first event occurs after 30 minutes.
	\end{enumerate}
\end{solucion}

\begin{solucion}
	\label{sol:4.5.4}
	For Problem \ref{prob:4.5.4}:
	\begin{enumerate}
		\item Proof of the memorylessness property:
		\begin{align*}
			P(X > s + t \mid X > s) = \dfrac{P(X > s + t)}{P(X > s)} = \dfrac{e^{-\lambda(s+t)}}{e^{-\lambda s}} = e^{-\lambda t} = P(X > t).
		\end{align*}
		\item For uniqueness: let $g(x) = P(X > x)$. The property implies $g(s + t) = g(s) g(t)$ for $s, t > 0$. The only measurable functions satisfying this Cauchy functional equation with $g(0) = 1$ and $g$ decreasing are $g(x) = e^{-cx}$ for some constant $c > 0$. Therefore, the only continuous distribution with memorylessness and support in $[0, \infty)$ is the exponential. \quad \qedhere
	\end{enumerate}
\end{solucion}

\begin{solucion}
	\label{sol:4.5.5}
	For Problem \ref{prob:4.5.5} with $X \sim \text{Exp}(0.5)$: the CDF is $F(x) = 1 - e^{-0.5 x}$, so $q_{p} = -2\ln(1-p)$.
	\begin{align*}
		q_{0.10} = -2\ln(0.9) \approx 0.211 \text{ min}, \quad q_{0.50} = 2\ln 2 \approx 1.386 \text{ min}, \quad q_{0.90} = -2\ln(0.1) \approx 4.605 \text{ min}.
	\end{align*}
	The quantile $q_{0.95} = -2\ln(0.05) \approx 5.991$ min means that 95\% of customers experience a waiting time shorter than 6 minutes. This quantile is fundamental for service design: it guarantees that the queue does not exceed a certain threshold.
\end{solucion}

\begin{solucion}
	\label{sol:4.5.6}
	For Problem \ref{prob:4.5.6} with `np.random.seed(42)` and `np.random.exponential(0.5, 100000)`:
	\begin{align*}
		\bar{X} \approx 0.5000 \quad (\text{theoretical: } 1/\lambda = 0.5), \quad s^{2} \approx 0.2500 \quad (\text{theoretical: } 1/\lambda^{2} = 0.25).
	\end{align*}
	The histogram shows a distribution with maximum density at $x = 0$ ($f(0) = \lambda = 2$) decaying exponentially, confirming the theoretical PDF.
\end{solucion}

\begin{solucion}
	\label{sol:4.5.7}
	For Problem \ref{prob:4.5.7}:
	\begin{enumerate}
		\item The log-likelihood for $n$ i.i.d. samples from $\text{Exp}(\lambda)$ is:
		\begin{align*}
			\ell(\lambda) = \sum_{i=1}^{n} \log f(x_i \mid \lambda) = n \log \lambda - \lambda \sum_{i=1}^{n} x_i.
		\end{align*}
		\item Differentiating and setting to zero: $\ell'(\lambda) = n/\lambda - \sum x_i = 0 \Rightarrow \hat{\lambda} = \dfrac{n}{\sum x_i} = \dfrac{1}{\bar{X}}$.
		\item Bias: $\E[\hat{\lambda}] = n \E[1/S]$ where $S = \sum X_i \sim \text{Gamma}(n, \lambda)$. Using $\E[1/S] = \lambda/(n-1)$ for $n > 1$, we obtain $\E[\hat{\lambda}] = n\lambda/(n-1) = \lambda \cdot n/(n-1)$, so $\hat{\lambda}$ has positive bias $\lambda/(n-1)$. However, $\hat{\lambda}_n \to \lambda$ as $n \to \infty$ (consistent). The asymptotic variance is $\Var(\hat{\lambda}) \to \lambda^{2}/n$ by the Fisher information $I(\lambda) = 1/\lambda^{2}$.
	\end{enumerate}
\end{solucion}

\begin{solucion}
	\label{sol:4.5.8}
	For Problem \ref{prob:4.5.8}:
	\begin{enumerate}
		\item The MGF of the sum is $M_{S_n}(t) = \prod_{i=1}^{n} M_{X_i}(t) = (\lambda/(\lambda - t))^{n}$ for $t < \lambda$ by independence. This is exactly the MGF of $\text{Gamma}(n, \lambda)$ (with integer shape parameter $n$, called the \emph{Erlang distribution}). By the uniqueness of the MGF, $S_n \sim \text{Gamma}(n, \lambda)$.
		\item For $\text{Gamma}(n, \lambda)$: $\E[S_n] = n/\lambda$ and $\Var(S_n) = n/\lambda^{2}$.
		\item For $n = 5$ and $\lambda = 1$: $S_5 \sim \text{Gamma}(5, 1)$ with $\E[S_5] = 5$ and $\sigma = \sqrt{5} \approx 2.236$. The PDF has a maximum at $x = 4$ (mode = $n - 1$ for the Gamma). By the CLT, $S_5 \approx N(5, 5)$, but the Gamma distribution is skewed for small $n$.
	\end{enumerate}
\end{solucion}

\begin{solucion}
	\label{sol:4.5.9}
	For Problem \ref{prob:4.5.9}:
	\begin{enumerate}
		\item By independence and the memorylessness property:
		\begin{align*}
			P(T > t) = P(\min(X_1, \dots, X_n) > t) = \prod_{i=1}^{n} P(X_i > t) = \prod_{i=1}^{n} e^{-\lambda t} = e^{-n\lambda t}.
		\end{align*}
		This is the tail of $\text{Exp}(n\lambda)$, so $T \sim \text{Exp}(n\lambda)$.
		\item $\E[T] = 1/(n\lambda)$. With $n = 10$ and $\lambda = 0.001$: $\E[T] = 100$ hours $\approx 4.17$ days.
		\item $P(T > 100) = e^{-10 \cdot 0.001 \cdot 100} = e^{-1} \approx 0.368$. The median is $m = \ln 2 / (n\lambda) = 0.693/0.01 = 69.3$ hours $\approx 2.89$ days. The system has a 36.8\% probability of surviving more than 100 hours, with a mean time to failure of 4.17 days.
	\end{enumerate}
\end{solucion}

\begin{solucion}
	\label{sol:4.5.10}
	For Problem \ref{prob:4.5.10} ($M/M/1$ queueing system):
	\begin{enumerate}
		\item In a Poisson process with rate $\lambda$, the times between consecutive arrivals are i.i.d. $\text{Exp}(\lambda)$ (fundamental property of the Poisson process). Analogously, if the service times are i.i.d. $\text{Exp}(\mu)$, the server attends each customer in a random time with mean $1/\mu$.
		\item The traffic intensity is $\rho = \lambda/\mu$, defined as the ratio between the arrival rate and the service rate. The system is stable if $\rho < 1$ (the service rate exceeds the arrival rate). In steady state, the probability that the system is busy (server active) is exactly $\rho$.
		\item With $\lambda = 5$ and $\mu = 6$: $\rho = 5/6 \approx 0.833$. The server is busy 83.3\% of the time in steady state, and there is a 16.7\% probability that it is idle. Since $\rho < 1$, the system is stable but with high utilization. \quad \qedhere
	\end{enumerate}
\end{solucion}

% ==============================================================================
% SECTION 04.06: NORMAL / GAUSSIAN DISTRIBUTION AND THE Z-SCORE
% ==============================================================================

\subsection*{Problem Statements --- Section 04.06}

\subsubsection*{Fundamental Level}

\begin{problema}
	\label{prob:4.6.1}
	A continuous random variable $X$ has a Normal distribution with mean $\mu = 100$ and standard deviation $\sigma = 15$.
	\begin{enumerate}
		\item Write the PDF $f(x) = \frac{1}{\sigma\sqrt{2\pi}} e^{-(x-\mu)^{2}/(2\sigma^{2})}$ and identify its components.
		\item Compute the probability $P(85 \le X \le 115)$ using the standardization $Z = (X - \mu)/\sigma$.
		\item Compute $P(X > 130)$ and $P(X < 70)$.
	\end{enumerate}
\end{problema}

\begin{problema}
	\label{prob:4.6.2}
	Let $X \sim N(\mu, \sigma^{2})$. Prove that the standardization $Z = (X - \mu)/\sigma$ follows the standard Normal $N(0, 1)$ with PDF $\phi(z) = \frac{1}{\sqrt{2\pi}} e^{-z^{2}/2}$.
	\begin{enumerate}
		\item Compute the expectation $\E[Z] = 0$ and the variance $\Var(Z) = 1$.
		\item Verify the symmetry: $\phi(-z) = \phi(z)$.
	\end{enumerate}
\end{problema}

\begin{problema}
	\label{prob:4.6.3}
	In a manufacturing process, the length of a metal shaft follows $N(10.0, 0.04)$ (in cm, with $\sigma = 0.2$ cm). The technical specifications require the length to be in $[9.7, 10.3]$ cm.
	\begin{enumerate}
		\item Compute the probability that a shaft meets the specifications using the $Z$-score.
		\item Determine the quantiles $z_{0.05}$ and $z_{0.95}$ such that $P(Z \le z_{0.05}) = 0.05$ and $P(Z \le z_{0.95}) = 0.95$.
		\item Compute the 90\% tolerance interval: $[L, U]$ such that $P(L \le X \le U) = 0.90$.
	\end{enumerate}
\end{problema}

\subsubsection*{Operational Level}

\begin{problema}
	\label{prob:4.6.4}
	The height of an adult population follows $N(170, 100)$ (in cm, $\sigma = 10$ cm).
	\begin{enumerate}
		\item Compute the probability that a randomly chosen person is taller than 185 cm.
		\item Compute the probability that a person is between 160 cm and 180 cm tall.
		\item How many people out of every 1000 are expected to be shorter than 155 cm?
	\end{enumerate}
\end{problema}

\begin{problema}
	\label{prob:4.6.5}
	Implement a Monte Carlo simulation with $N = 250{,}000$ samples from $N(100, 225)$ ($\sigma = 15$) and verify:
	\begin{enumerate}
		\item The empirical mean and standard deviation match $\mu = 100$ and $\sigma = 15$.
		\item The empirical rule: $P(\mu - \sigma \le X \le \mu + \sigma) \approx 0.6827$.
		\item The skewness and kurtosis of the samples are close to 0 and 0 (characteristics of the Normal).
	\end{enumerate}
\end{problema}

\begin{problema}
	\label{prob:4.6.6}
	In hypothesis testing, the statistic $Z = (\bar{X} - \mu_{0})/(\sigma/\sqrt{n})$ follows $N(0, 1)$ under $H_{0}: \mu = \mu_{0}$. For a sample of $n = 25$ observations with $\sigma = 10$ and $\bar{x} = 103$:
	\begin{enumerate}
		\item Compute the statistic $Z$ under $H_{0}: \mu = 100$.
		\item Compute the two-tailed $p$-value: $p = 2(1 - \Phi(|Z|))$.
		\item Decide whether we reject $H_{0}$ at level $\alpha = 0.05$.
	\end{enumerate}
\end{problema}

\subsubsection*{Analytical Level}

\begin{problema}
	\label{prob:4.6.7}
	Prove that if $X \sim N(\mu, \sigma^{2})$, its moment generating function is $M_{X}(t) = e^{\mu t + \sigma^{2} t^{2}/2}$.
	\begin{enumerate}
		\item Complete the square in the exponent of the integral $M_{X}(t) = \int e^{tx} f(x)\,dx$.
		\item Use the Gaussian integral $\int e^{-u^{2}/2}\,du = \sqrt{2\pi}$.
		\item Derive the moments: $\E[X] = M'(0) = \mu$ and $\E[X^{2}] = M''(0) = \mu^{2} + \sigma^{2}$.
	\end{enumerate}
\end{problema}

\begin{problema}
	\label{prob:4.6.8}
	Let $X_{1}, \dots, X_{n}$ be independent i.i.d. $N(\mu, \sigma^{2})$ variables. Prove that the sample mean $\bar{X} = (1/n) \sum X_{i}$ follows $N(\mu, \sigma^{2}/n)$, and that the sample variance $S^{2} = \frac{1}{n-1}\sum(X_{i} - \bar{X})^{2}$ satisfies $(n-1)S^{2}/\sigma^{2} \sim \chi^{2}_{n-1}$.
\end{problema}

\subsubsection*{Challenging Level}

\begin{problema}
	\label{prob:4.6.9}
	State and prove the \emph{Central Limit Theorem}: if $X_{1}, \dots, X_{n}$ are i.i.d. with mean $\mu$ and variance $\sigma^{2} < \infty$, then $Z_{n} = (\bar{X} - \mu)/(\sigma/\sqrt{n}) \xrightarrow{d} N(0, 1)$ as $n \to \infty$. Hint: use the moment generating function and expand it in a Taylor series.
\end{problema}

\begin{problema}
	\label{prob:4.6.10}
	On a standardized exam with $n = 400$ multiple-choice questions with 4 options each, the number of correct answers $X$ follows $X \sim \text{Bin}(400, 0.25)$. Use the Normal approximation with the Yates continuity correction to compute:
	\begin{enumerate}
		\item $P(X \ge 120)$ using $X \approx N(100, 75)$.
		\item Compare with the exact value from the Binomial.
		\item Determine the $n$ needed so that $\text{SE}(\hat{p}) = \sqrt{p(1-p)/n} < 0.01$ with $p = 0.25$.
	\end{enumerate}
\end{problema}

\subsection*{Hints for the Problems --- Section 04.06}

\begin{sugerencia}
	\label{sug:4.6.1}
	For Problem \ref{prob:4.6.1}, $f(x) = \frac{1}{15\sqrt{2\pi}} e^{-(x-100)^{2}/450}$. We standardize: $Z = (X-100)/15$. $P(85 \le X \le 115) = P(-1 \le Z \le 1) = 2\Phi(1) - 1 \approx 0.6827$. $P(X > 130) = P(Z > 2) = 1 - \Phi(2) \approx 0.0228$. $P(X < 70) = P(Z < -2) = \Phi(-2) \approx 0.0228$.
\end{sugerencia}

\begin{sugerencia}
	\label{sug:4.6.2}
	For Problem \ref{prob:4.6.2}, use the change of variable $z = (x - \mu)/\sigma$ in the PDF integral: $f_{Z}(z) = f_{X}(\sigma z + \mu) \cdot \sigma = \frac{1}{\sqrt{2\pi}} e^{-z^{2}/2}$. $\E[Z] = \int z \phi(z)\,dz = 0$ by symmetry. $\Var(Z) = \E[Z^{2}] = 1$ by direct computation.
\end{sugerencia}

\begin{sugerencia}
	\label{sug:4.6.3}
	For Problem \ref{prob:4.6.3}, $P(9.7 \le X \le 10.3) = P(-1.5 \le Z \le 1.5) = 2\Phi(1.5) - 1 \approx 0.8664$. The typical quantiles: $z_{0.05} \approx -1.645$ and $z_{0.95} \approx 1.645$. The 90\% tolerance interval is $[10 - 1.645 \cdot 0.2, 10 + 1.645 \cdot 0.2] = [9.671, 10.329]$.
\end{sugerencia}

\begin{sugerencia}
	\label{sug:4.6.4}
	For Problem \ref{prob:4.6.4}, $P(X > 185) = P(Z > 1.5) = 1 - \Phi(1.5) \approx 0.0668$ (6.68\% of the population). $P(160 \le X \le 180) = P(-1 \le Z \le 1) \approx 0.6827$. We expect $0.1586 \cdot 1000 \approx 159$ people with height $< 155$ cm.
\end{sugerencia}

\begin{sugerencia}
	\label{sug:4.6.5}
	For Problem \ref{prob:4.6.5}, use `np.random.seed(42)` and `np.random.normal(100, 15, 250000)`. Verify `np.mean(samples)`, `np.std(samples, ddof=1)`, and compute `np.mean((samples - 100)**3 / 15**3)` for the skewness. The 68-95-99.7 rule: $\approx 68.27\%$ within $\pm 1\sigma$, $\approx 95.45\%$ within $\pm 2\sigma$.
\end{sugerencia}

\begin{sugerencia}
	\label{sug:4.6.6}
	For Problem \ref{prob:4.6.6}, $Z = (103 - 100)/(10/\sqrt{25}) = 3/2 = 1.5$. The two-tailed $p$-value is $p = 2(1 - \Phi(1.5)) \approx 2(0.0668) = 0.1336$. Since $p = 0.1336 > 0.05$, we do not reject $H_{0}$ at level $\alpha = 0.05$.
\end{sugerencia}

\begin{sugerencia}
	\label{sug:4.6.7}
	For Problem \ref{prob:4.6.7}, complete the square: $tx - (x-\mu)^{2}/(2\sigma^{2}) = -(x - (\mu + \sigma^{2}t))^{2}/(2\sigma^{2}) + \mu t + \sigma^{2}t^{2}/2$. Factoring $e^{\mu t + \sigma^{2}t^{2}/2}$ out of the integral, what remains is the Normal PDF evaluated at $\mu + \sigma^{2}t$, which integrates to 1.
\end{sugerencia}

\begin{sugerencia}
	\label{sug:4.6.8}
	For Problem \ref{prob:4.6.8}, $\bar{X} = (1/n) \sum X_{i}$ has MGF $M_{\bar{X}}(t) = \prod e^{\mu t/n + \sigma^{2}t^{2}/(2n^{2})} = e^{\mu t + \sigma^{2}t^{2}/(2n)}$, confirming $N(\mu, \sigma^{2}/n)$. For $S^{2}$: Fisher's theorem states that $(n-1)S^{2}/\sigma^{2} \sim \chi^{2}_{n-1}$ independently of $\bar{X}$ (Cochran's theorem).
\end{sugerencia}

\begin{sugerencia}
	\label{sug:4.6.9}
	For Problem \ref{prob:4.6.9}, the MGF of $Z_{n}$ is $M_{Z_n}(t) = e^{-\mu t/\sigma_n} M_{X_{1}}(t/\sigma_n)^{n}$. Applying the Taylor expansion of $\log M_{X_{1}}$ around 0: $\log M_{X_{1}}(t/\sigma_n) = \mu t/\sigma_n + \sigma^{2}t^{2}/(2n) + O(1/n^{3/2})$. Therefore, $M_{Z_{n}}(t) \to e^{t^{2}/2}$ as $n \to \infty$, which is the MGF of $N(0, 1)$.
\end{sugerencia}

\begin{sugerencia}
	\label{sug:4.6.10}
	For Problem \ref{prob:4.6.10}, $\mu = np = 100$ and $\sigma = \sqrt{np(1-p)} = \sqrt{75} \approx 8.66$. With the Yates correction: $P(X \ge 120) = P(Y \ge 119.5)$ with $Y \sim N(100, 75)$: $Z = (119.5 - 100)/8.66 \approx 2.25$, $P(Z \ge 2.25) \approx 0.0122$. The exact value is $1 - \text{binom.cdf}(119, 400, 0.25) \approx 0.0098$. For $\text{SE} < 0.01$: $n > 0.25 \cdot 0.75/0.0001 = 1875$.
\end{sugerencia}

\subsection*{Solutions to the Problems --- Section 04.06}

\begin{solucion}
	\label{sol:4.6.1}
	For Problem \ref{prob:4.6.1}:
	\begin{enumerate}
		\item $f(x) = \dfrac{1}{15\sqrt{2\pi}} \exp\left(-\dfrac{(x-100)^{2}}{450}\right)$. The factor $1/(15\sqrt{2\pi}) \approx 0.0266$ normalizes the integral to 1; the exponential term concentrates the mass around $x = 100$.
		\item $P(85 \le X \le 115) = P(-1 \le Z \le 1) = 2\Phi(1) - 1 \approx 0.6827$.
		\item $P(X > 130) = P(Z > 2) = 1 - \Phi(2) \approx 1 - 0.9772 = 0.0228$. $P(X < 70) = P(Z < -2) = \Phi(-2) \approx 0.0228$.
	\end{enumerate}
\end{solucion}

\begin{solucion}
	\label{sol:4.6.2}
	For Problem \ref{prob:4.6.2}: with the change of variable $z = (x - \mu)/\sigma$, $dx = \sigma\,dz$, and support $\mathbb{R} \to \mathbb{R}$:
	\begin{align*}
		f_{Z}(z) = f_{X}(\sigma z + \mu) \cdot \sigma = \dfrac{1}{\sigma\sqrt{2\pi}} e^{-(\sigma z)^{2}/(2\sigma^{2})} \cdot \sigma = \dfrac{1}{\sqrt{2\pi}} e^{-z^{2}/2}.
	\end{align*}
	This is the standard Normal PDF $\phi(z)$. $\E[Z] = 0$ and $\Var(Z) = 1$ by direct computation: $\int z\phi(z)\,dz = 0$ by symmetry, and $\int z^{2}\phi(z)\,dz = 1$. \quad \qedhere
\end{solucion}

\begin{solucion}
	\label{sol:4.6.3}
	For Problem \ref{prob:4.6.3} with $X \sim N(10, 0.04)$:
	\begin{enumerate}
		\item $P(9.7 \le X \le 10.3) = P(-1.5 \le Z \le 1.5) = 2\Phi(1.5) - 1 \approx 0.8664$. Approximately 86.6\% of shafts meet the specifications.
		\item The typical quantiles are $z_{0.05} \approx -1.645$ and $z_{0.95} \approx 1.645$.
		\item The 90\% tolerance interval is $[10 - 1.645 \cdot 0.2, 10 + 1.645 \cdot 0.2] = [9.671, 10.329]$ cm. 90\% of shafts fall in this interval.
	\end{enumerate}
\end{solucion}

\begin{solucion}
	\label{sol:4.6.4}
	For Problem \ref{prob:4.6.4} with $X \sim N(170, 100)$ ($\sigma = 10$):
	\begin{enumerate}
		\item $P(X > 185) = P(Z > 1.5) = 1 - \Phi(1.5) \approx 0.0668$. Approximately 6.68\% of the population is taller than 185 cm.
		\item $P(160 \le X \le 180) = P(-1 \le Z \le 1) \approx 0.6827$. 68.27\% of the population is between 160 and 180 cm tall.
		\item $P(X < 155) = P(Z < -1.5) \approx 0.0668$. We expect $0.0668 \cdot 1000 \approx 67$ people with height less than 155 cm.
	\end{enumerate}
\end{solucion}

\begin{solucion}
	\label{sol:4.6.5}
	For Problem \ref{prob:4.6.5} with `np.random.seed(42)` and `np.random.normal(100, 15, 250000)`:
	\begin{itemize}
		\item $\bar{X} \approx 100.0$ (theoretical: 100), $s \approx 15.0$ (theoretical: 15).
		\item $P(85 \le X \le 115) \approx 0.6827$ (68-95-99.7 rule).
		\item Empirical skewness $\approx 0$, kurtosis $\approx 0$, confirming the symmetry of the Normal.
	\end{itemize}
\end{solucion}

\begin{solucion}
	\label{sol:4.6.6}
	For Problem \ref{prob:4.6.6}:
	\begin{enumerate}
		\item $Z = (103 - 100)/(10/\sqrt{25}) = 3/2 = 1.5$.
		\item The two-tailed $p$-value is $p = 2(1 - \Phi(1.5)) \approx 2(0.0668) = 0.1336$.
		\item Since $p = 0.1336 > 0.05$, we do not reject $H_{0}$ at level $\alpha = 0.05$. The sample does not provide significant evidence that $\mu \ne 100$.
	\end{enumerate}
\end{solucion}

\begin{solucion}
	\label{sol:4.6.7}
	For Problem \ref{prob:4.6.7}: the MGF is
	\begin{align*}
		M_{X}(t) = \int_{-\infty}^{\infty} e^{tx} \cdot \dfrac{1}{\sigma\sqrt{2\pi}} e^{-(x-\mu)^{2}/(2\sigma^{2})}\,dx.
	\end{align*}
	Completing the square in the exponent: $tx - (x-\mu)^{2}/(2\sigma^{2}) = -(x - (\mu + \sigma^{2}t))^{2}/(2\sigma^{2}) + \mu t + \sigma^{2}t^{2}/2$. Factoring:
	\begin{align*}
		M_{X}(t) = e^{\mu t + \sigma^{2}t^{2}/2} \cdot \int \dfrac{1}{\sigma\sqrt{2\pi}} e^{-(x - (\mu + \sigma^{2}t))^{2}/(2\sigma^{2})}\,dx = e^{\mu t + \sigma^{2}t^{2}/2},
	\end{align*}
	where the integral is exactly 1 (it is the Normal PDF with mean $\mu + \sigma^{2}t$ and variance $\sigma^{2}$). Differentiating: $\E[X] = M'(0) = \mu$ and $\E[X^{2}] = M''(0) = \mu^{2} + \sigma^{2}$. \quad \qedhere
\end{solucion}

\begin{solucion}
	\label{sol:4.6.8}
	For Problem \ref{prob:4.6.8}: the MGF of $\bar{X}$ is
	\begin{align*}
		M_{\bar{X}}(t) = \E\left[e^{t\bar{X}}\right] = \E\left[e^{(t/n)\sum X_{i}}\right] = \prod_{i=1}^{n} e^{\mu t/n + \sigma^{2}t^{2}/(2n^{2})} = e^{\mu t + \sigma^{2}t^{2}/(2n)},
	\end{align*}
	which is the MGF of $N(\mu, \sigma^{2}/n)$. For $S^{2}$: Fisher's theorem establishes that $(n-1)S^{2}/\sigma^{2} = \sum(X_{i} - \bar{X})^{2}/\sigma^{2} \sim \chi^{2}_{n-1}$, independently of $\bar{X}$ (Cochran). This independence is fundamental for the construction of the Student's $t$ distribution. \quad \qedhere
\end{solucion}

\begin{solucion}
	\label{sol:4.6.9}
	For Problem \ref{prob:4.6.9} (CLT), the MGF of $Z_{n} = (\bar{X} - \mu)/(\sigma/\sqrt{n})$ is
	\begin{align*}
		M_{Z_{n}}(t) = e^{-\mu t/\sigma_n} \cdot M_{X_{1}}(t/\sigma_n)^{n}, \quad \sigma_n = \sigma/\sqrt{n}.
	\end{align*}
	Applying the Taylor expansion of $\log M_{X_{1}}$ around 0: $\log M_{X_{1}}(s) = \mu s + \sigma^{2}s^{2}/2 + O(s^{3})$. With $s = t/\sigma_n$:
	\begin{align*}
		\log M_{X_{1}}(t/\sigma_n) = \dfrac{\mu t}{\sigma_n} + \dfrac{\sigma^{2}t^{2}}{2n} + O(n^{-3/2}).
	\end{align*}
	Multiplying by $n$ and subtracting $n\mu t/\sigma_n$: $\log M_{Z_{n}}(t) = t^{2}/2 + O(n^{-1/2}) \to t^{2}/2$. By the uniqueness of the MGF and Lévy's Continuity Theorem, $M_{Z_{n}}(t) \to e^{t^{2}/2}$, which is the MGF of $N(0, 1)$. \quad \qedhere
\end{solucion}

\begin{solucion}
	\label{sol:4.6.10}
	For Problem \ref{prob:4.6.10} with $n = 400$, $p = 0.25$:
	\begin{enumerate}
		\item $\mu = np = 100$ and $\sigma = \sqrt{np(1-p)} = \sqrt{75} \approx 8.66$. With the Yates correction: $P(X \ge 120) \approx P(Y \ge 119.5) = P(Z \ge (119.5 - 100)/8.66) = P(Z \ge 2.25) \approx 0.0122$.
		\item The exact value is $P(X \ge 120) = 1 - \text{binom.cdf}(119, 400, 0.25) \approx 0.0133$. The absolute error is $\approx 0.0011$, confirming the accuracy of the Yates correction.
		\item For $\text{SE}(\hat{p}) = \sqrt{p(1-p)/n} < 0.01$: $n > p(1-p)/0.0001 = 0.1875/0.0001 = 1875$. At least $n = 1875$ observations are needed.
	\end{enumerate}
\end{solucion}

% ==============================================================================
% SECTION 04.07: GAMMA, BETA, AND WEIBULL DISTRIBUTIONS
% ==============================================================================

\subsection*{Problem Statements --- Section 04.07}

\subsubsection*{Fundamental Level}

\begin{problema}
	\label{prob:4.7.1}
	Let $X \sim \text{Gamma}(\alpha=3, \beta=2)$.
	\begin{enumerate}
		\item Write the PDF $f(x)$ and identify the shape parameter and the scale parameter.
		\item Compute the mean $\mu_X = \alpha\beta$ and the variance $\sigma_X^{2} = \alpha\beta^{2}$.
		\item Verify, using the definition of the gamma function $\Gamma(\alpha) = \int_0^\infty t^{\alpha-1}e^{-t}\,dt$, that $\int_0^\infty f(x)\,dx = 1$.
	\end{enumerate}
\end{problema}

\begin{problema}
	\label{prob:4.7.2}
	Identify the following special cases of the gamma family.
	\begin{enumerate}
		\item Let $X \sim \text{Gamma}(1, 4)$. Identify the equivalent distribution of $X$ and compute $P(X > 6)$.
		\item Let $Y \sim \text{Gamma}(5/2, 2)$. Identify $Y$ as a chi-squared distribution and state its degrees of freedom.
	\end{enumerate}
\end{problema}

\begin{problema}
	\label{prob:4.7.3}
	The moment generating function of $X \sim \text{Gamma}(\alpha, \beta)$ is $M_{X}(t) = (1-\beta t)^{-\alpha}$, for $t < 1/\beta$.
	\begin{enumerate}
		\item Compute $M_{X}'(0)$ to obtain $\E[X] = \alpha\beta$.
		\item Compute $M_{X}''(0)$ to obtain $\E[X^{2}]$ and from there $\Var(X) = \alpha\beta^{2}$.
	\end{enumerate}
\end{problema}

\subsubsection*{Operational Level}

\begin{problema}
	\label{prob:4.7.4}
	Let $X \sim \text{Beta}(\alpha=3, \beta=5)$.
	\begin{enumerate}
		\item Write the PDF $f(x)$ using the beta function $B(\alpha,\beta)$.
		\item Compute $B(3,5)$ using its relation to the gamma function, $B(\alpha,\beta) = \Gamma(\alpha)\Gamma(\beta)/\Gamma(\alpha+\beta)$.
		\item Compute $\E(X)$ and $\Var(X)$.
	\end{enumerate}
\end{problema}

\begin{problema}
	\label{prob:4.7.5}
	In a satisfaction survey, the proportion $X$ of customers satisfied with a new product is modeled as $X \sim \text{Beta}(6,4)$ according to historical data.
	\begin{enumerate}
		\item Compute the expected proportion of satisfied customers.
		\item Compute the variance and standard deviation of $X$.
		\item Compare the mean with the mode $(\alpha-1)/(\alpha+\beta-2)$ and determine whether $X$ is more likely to be above or below $0.5$.
	\end{enumerate}
\end{problema}

\begin{problema}
	\label{prob:4.7.6}
	Let $X \sim \text{Beta}(\alpha,\beta)$ be the prior on the probability of success $p$ of a Bernoulli process, with initial prior $\text{Beta}(2,2)$. We observe $n=20$ trials with $k=14$ successes.
	\begin{enumerate}
		\item Determine the posterior distribution of $p$ (recall that the conjugate posterior is $\text{Beta}(\alpha+k, \beta+n-k)$).
		\item Compute the posterior mean $\E(p \mid \text{data})$.
		\item State the symmetry property $\text{Beta}(\alpha,\beta) = 1 - \text{Beta}(\beta,\alpha)$ and explain, using the change of variable $u = 1-x$ in the PDF integral, why it holds.
	\end{enumerate}
\end{problema}

\subsubsection*{Analytical Level}

\begin{problema}
	\label{prob:4.7.7}
	The lifetime $T$ (in years) of a mechanical component follows $T \sim \text{Weibull}(\beta=1.5, \eta=4)$.
	\begin{enumerate}
		\item Write the PDF $f(t)$ and the reliability function $R(t) = 1 - F(t)$.
		\item Compute $P(T > 3)$.
		\item Determine whether the failure rate $h(t)$ is increasing or decreasing and explain its physical meaning.
	\end{enumerate}
\end{problema}

\begin{problema}
	\label{prob:4.7.8}
	Two batches of components have Weibull lifetimes with the same scale $\eta = 10$: Batch A $\sim \text{Weibull}(\beta=1, \eta=10)$ and Batch B $\sim \text{Weibull}(\beta=2, \eta=10)$.
	\begin{enumerate}
		\item Compute the mean of each batch using $\mu = \eta\,\Gamma(1+1/\beta)$ (recall $\Gamma(2)=1$ and $\Gamma(1.5) = \sqrt{\pi}/2 \approx 0.8862$).
		\item Compare the hazard functions $h_{A}(t)$ and $h_{B}(t)$ at $t=5$ and $t=15$. Which batch is more reliable in the short term? And in the long term?
		\item Explain the physical difference between the two batches in terms of component aging.
	\end{enumerate}
\end{problema}

\subsubsection*{Challenging Level}

\begin{problema}
	\label{prob:4.7.9}
	Let $X \sim \text{Gamma}(\alpha,\beta)$ with $\alpha = \nu/2$ and $\beta = 2$ (so that $X \sim \chi^{2}_{\nu}$).
	\begin{enumerate}
		\item Use the MGF of the Gamma to show that $M_{X}(t) = (1-2t)^{-\nu/2}$, which coincides with the MGF of the chi-squared distribution.
		\item Knowing that $Z^{2} \sim \text{Gamma}(1/2, 2) = \chi^{2}_{1}$ for $Z \sim N(0,1)$, use the property of the sum of independent gamma variables to justify that the sum of $\nu$ independent $\text{Gamma}(1/2,2)$ variables is $\text{Gamma}(\nu/2, 2)$.
		\item Apply the above to $S = Z_{1}^{2}+Z_{2}^{2}+Z_{3}^{2}$ with $Z_{i} \sim N(0,1)$ i.i.d., and compute $P(S > 7.815)$ using the fact that $7.815$ is the $0.95$ quantile of $\chi^{2}_{3}$.
	\end{enumerate}
\end{problema}

\begin{problema}
	\label{prob:4.7.10}
	A spare-parts system replaces a critical component as soon as it fails. The lifetimes between replacements are i.i.d. $\text{Exp}(\lambda=0.5)$ (failures per year, equivalent to $\text{Gamma}(1,2)$).
	\begin{enumerate}
		\item What is the distribution of the time $T$ until the fifth replacement? Use the additive property of the Gamma.
		\item Compute $\E(T)$ and $\Var(T)$.
		\item Compare this additive model (Gamma/Erlang process) with a $\text{Weibull}(\beta=1,\eta=2)$ model for the lifetime of a single component, explaining why both coincide when $\beta_{\text{Weibull}}=1$.
	\end{enumerate}
\end{problema}

\subsection*{Hints for the Problems --- Section 04.07}

\begin{sugerencia}
	\label{sug:4.7.1}
	For Problem \ref{prob:4.7.1}, $f(x) = \frac{1}{\Gamma(3)\cdot 2^{3}}x^{2}e^{-x/2} = \frac{x^{2}e^{-x/2}}{16}$ for $x>0$, with $\Gamma(3)=2!=2$. $\mu_X = 3\cdot 2=6$, $\sigma_X^{2}=3\cdot 4=12$. For the normalization, substitute $u=x/2$: $\int_0^\infty f(x)\,dx = \frac{1}{\Gamma(3)}\int_0^\infty u^{2}e^{-u}\,du = \frac{\Gamma(3)}{\Gamma(3)}=1$.
\end{sugerencia}

\begin{sugerencia}
	\label{sug:4.7.2}
	For Problem \ref{prob:4.7.2}, $\text{Gamma}(1,\beta) = \text{Exp}(\lambda=1/\beta)$, so $X \sim \text{Exp}(0.25)$ and $P(X>6)=e^{-0.25\cdot 6}=e^{-1.5}$. For $Y$, use $\alpha=\nu/2 \implies \nu=2\alpha=5$.
\end{sugerencia}

\begin{sugerencia}
	\label{sug:4.7.3}
	For Problem \ref{prob:4.7.3}, differentiate $M_{X}(t)=(1-\beta t)^{-\alpha}$ term by term: $M_X'(t)=\alpha\beta(1-\beta t)^{-\alpha-1}$ and $M_X''(t)=\alpha(\alpha+1)\beta^{2}(1-\beta t)^{-\alpha-2}$. Evaluate at $t=0$ and use $\Var(X)=M_X''(0)-(M_X'(0))^{2}$.
\end{sugerencia}

\begin{sugerencia}
	\label{sug:4.7.4}
	For Problem \ref{prob:4.7.4}, $f(x) = x^{2}(1-x)^{4}/B(3,5)$ on $(0,1)$. Since $\Gamma(3)=2$, $\Gamma(5)=24$, and $\Gamma(8)=5040$, $B(3,5) = 2\cdot 24/5040 = 1/105$. $\E(X)=3/8=0.375$, $\Var(X)=15/(64\cdot 9)=15/576\approx 0.0260$.
\end{sugerencia}

\begin{sugerencia}
	\label{sug:4.7.5}
	For Problem \ref{prob:4.7.5}, $\E(X)=6/10=0.6$, $\Var(X)=24/(100\cdot 11)=24/1100\approx 0.0218$, $\text{sd}(X)\approx 0.148$. The mode is $5/8=0.625$, greater than the mean, indicating that the mass is more concentrated above $0.5$.
\end{sugerencia}

\begin{sugerencia}
	\label{sug:4.7.6}
	For Problem \ref{prob:4.7.6}, the posterior is $\text{Beta}(2+14,\,2+20-14)=\text{Beta}(16,8)$, with mean $16/24=2/3$. For the symmetry, substitute $u=1-x$, $du=-dx$: the integral of $\text{Beta}(\alpha,\beta)$ in $x$ transforms into the integral of $\text{Beta}(\beta,\alpha)$ in $u$ with the limits reversed.
\end{sugerencia}

\begin{sugerencia}
	\label{sug:4.7.7}
	For Problem \ref{prob:4.7.7}, $R(t)=\exp[-(t/4)^{1.5}]$. For $t=3$: $(3/4)^{1.5}=0.75\sqrt{0.75}\approx 0.6495$, so $P(T>3)=e^{-0.6495}\approx 0.522$. Since $\beta=1.5>1$, the failure rate is increasing (wear-out).
\end{sugerencia}

\begin{sugerencia}
	\label{sug:4.7.8}
	For Problem \ref{prob:4.7.8}, $\mu_A = 10\cdot\Gamma(2)=10$ and $\mu_B=10\cdot\Gamma(1.5)\approx 8.862$. $h_A(t)=1/10=0.1$ (constant). $h_B(t)=(2/10)(t/10)=t/50$: at $t=5$, $h_B=0.1$; at $t=15$, $h_B=0.3$.
\end{sugerencia}

\begin{sugerencia}
	\label{sug:4.7.9}
	For Problem \ref{prob:4.7.9}, substitute $\alpha=\nu/2,\beta=2$ in $M_X(t)=(1-\beta t)^{-\alpha}$. For part (c), $S\sim\chi^{2}_{3}$ and $7.815$ is by definition the $0.95$ quantile, so $P(S>7.815)=0.05$ directly from the definition of the quantile.
\end{sugerencia}

\begin{sugerencia}
	\label{sug:4.7.10}
	For Problem \ref{prob:4.7.10}, $T_i \sim \text{Gamma}(1,2)$ i.i.d.; by the additive property, $T=\sum_{i=1}^{5}T_i \sim \text{Gamma}(5,2)$. $\E(T)=5\cdot 2=10$, $\Var(T)=5\cdot 4=20$. $\text{Weibull}(1,\eta)$ coincides with $\text{Exp}(1/\eta)$, which is the distribution of each individual $T_i$, not of the sum $T$.
\end{sugerencia}

\subsection*{Solutions to the Problems --- Section 04.07}

\begin{solucion}
	\label{sol:4.7.1}
	For Problem \ref{prob:4.7.1} with $X \sim \text{Gamma}(3,2)$:
	\begin{enumerate}
		\item $f(x) = \dfrac{1}{\Gamma(3)\cdot 2^{3}}x^{2}e^{-x/2} = \dfrac{x^{2}e^{-x/2}}{16}$, $x>0$. The shape parameter is $\alpha=3$ and the scale parameter is $\beta=2$.
		\item $\mu_X = \alpha\beta = 3\cdot 2 = 6$. $\sigma_X^{2} = \alpha\beta^{2} = 3\cdot 4 = 12$.
		\item With the change of variable $u=x/2$ ($dx=2\,du$):
		\begin{align*}
			\int_0^\infty f(x)\,dx = \int_0^\infty \frac{x^{2}e^{-x/2}}{16}\,dx = \frac{1}{\Gamma(3)}\int_0^\infty u^{2}e^{-u}\,du = \frac{\Gamma(3)}{\Gamma(3)} = 1. \quad \qedhere
		\end{align*}
	\end{enumerate}
\end{solucion}

\begin{solucion}
	\label{sol:4.7.2}
	For Problem \ref{prob:4.7.2}:
	\begin{enumerate}
		\item $\text{Gamma}(1,4) = \text{Exp}(\lambda=1/4)$, so $X\sim\text{Exp}(0.25)$ and $P(X>6) = e^{-0.25\cdot 6} = e^{-1.5} \approx 0.2231$.
		\item Since $\alpha=\nu/2=5/2$, it follows that $\nu=5$; therefore $Y \sim \chi^{2}_{5}$. \quad \qedhere
	\end{enumerate}
\end{solucion}

\begin{solucion}
	\label{sol:4.7.3}
	For Problem \ref{prob:4.7.3}: differentiating $M_X(t)=(1-\beta t)^{-\alpha}$,
	\begin{align*}
		M_X'(t) = \alpha\beta(1-\beta t)^{-\alpha-1}, \qquad M_X''(t) = \alpha(\alpha+1)\beta^{2}(1-\beta t)^{-\alpha-2}.
	\end{align*}
	\begin{enumerate}
		\item $\E[X] = M_X'(0) = \alpha\beta$.
		\item $\E[X^{2}] = M_X''(0) = \alpha(\alpha+1)\beta^{2}$, so $\Var(X) = \alpha(\alpha+1)\beta^{2} - (\alpha\beta)^{2} = \alpha\beta^{2}$. \quad \qedhere
	\end{enumerate}
\end{solucion}

\begin{solucion}
	\label{sol:4.7.4}
	For Problem \ref{prob:4.7.4} with $X \sim \text{Beta}(3,5)$:
	\begin{enumerate}
		\item $f(x) = \dfrac{x^{2}(1-x)^{4}}{B(3,5)}$, $0<x<1$.
		\item $\Gamma(3)=2!=2$, $\Gamma(5)=4!=24$, $\Gamma(8)=7!=5040$, so $B(3,5) = \dfrac{2\cdot 24}{5040} = \dfrac{48}{5040} = \dfrac{1}{105} \approx 0.00952$.
		\item $\E(X) = \dfrac{3}{3+5} = 0.375$. $\Var(X) = \dfrac{3\cdot 5}{8^{2}\cdot 9} = \dfrac{15}{576} \approx 0.0260$. \quad \qedhere
	\end{enumerate}
\end{solucion}

\begin{solucion}
	\label{sol:4.7.5}
	For Problem \ref{prob:4.7.5} with $X \sim \text{Beta}(6,4)$:
	\begin{enumerate}
		\item $\E(X) = 6/10 = 0.6$, that is, $60\%$ of customers are expected to be satisfied.
		\item $\Var(X) = \dfrac{6\cdot 4}{10^{2}\cdot 11} = \dfrac{24}{1100} \approx 0.0218$, with standard deviation $\approx 0.1477$.
		\item The mode is $\dfrac{\alpha-1}{\alpha+\beta-2} = \dfrac{5}{8} = 0.625$, greater than the mean $0.6$. Since both the mean and the mode are above $0.5$, it is more likely that $X$ is above $0.5$. \quad \qedhere
	\end{enumerate}
\end{solucion}

\begin{solucion}
	\label{sol:4.7.6}
	For Problem \ref{prob:4.7.6}:
	\begin{enumerate}
		\item The conjugate posterior is $\text{Beta}(\alpha+k,\,\beta+n-k) = \text{Beta}(2+14,\,2+20-14) = \text{Beta}(16,8)$.
		\item $\E(p\mid\text{data}) = \dfrac{16}{16+8} = \dfrac{16}{24} = \dfrac{2}{3} \approx 0.667$.
		\item With $u=1-x$ (and $du=-dx$), the density of $\text{Beta}(\alpha,\beta)$ evaluated at $x$ becomes the density of $\text{Beta}(\beta,\alpha)$ evaluated at $u$:
		\begin{align*}
			f_X(x) = \frac{x^{\alpha-1}(1-x)^{\beta-1}}{B(\alpha,\beta)} \; \xrightarrow{u=1-x} \; \frac{(1-u)^{\alpha-1}u^{\beta-1}}{B(\alpha,\beta)} = \frac{u^{\beta-1}(1-u)^{\alpha-1}}{B(\beta,\alpha)},
		\end{align*}
		since $B(\alpha,\beta)=B(\beta,\alpha)$. Integrating both sides gives $P(X\le x) = P(1-X \ge 1-x)$, that is, $\text{Beta}(\alpha,\beta) = 1-\text{Beta}(\beta,\alpha)$. \quad \qedhere
	\end{enumerate}
\end{solucion}

\begin{solucion}
	\label{sol:4.7.7}
	For Problem \ref{prob:4.7.7} with $T \sim \text{Weibull}(1.5,4)$:
	\begin{enumerate}
		\item $f(t) = \dfrac{1.5}{4}\left(\dfrac{t}{4}\right)^{0.5}\exp\!\left[-\left(\dfrac{t}{4}\right)^{1.5}\right]$ and $R(t) = \exp\!\left[-\left(\dfrac{t}{4}\right)^{1.5}\right]$, for $t\ge 0$.
		\item $(3/4)^{1.5} = 0.75\cdot\sqrt{0.75} \approx 0.75\cdot 0.8660 \approx 0.6495$. Therefore $P(T>3) = R(3) = e^{-0.6495} \approx 0.522$.
		\item Since $\beta=1.5>1$, the failure rate $h(t) = \frac{1.5}{4}(t/4)^{0.5}$ is increasing: the component exhibits a wear-out pattern, with increasing failure risk as it ages. \quad \qedhere
	\end{enumerate}
\end{solucion}

\begin{solucion}
	\label{sol:4.7.8}
	For Problem \ref{prob:4.7.8}:
	\begin{enumerate}
		\item $\mu_A = 10\cdot\Gamma(1+1/1) = 10\cdot\Gamma(2) = 10\cdot 1 = 10$. $\mu_B = 10\cdot\Gamma(1+1/2) = 10\cdot\Gamma(1.5) \approx 10\cdot 0.8862 = 8.862$.
		\item $h_A(t) = \dfrac{1}{10}\left(\dfrac{t}{10}\right)^{0} = 0.1$ for all $t$ (constant rate, since $\beta=1$). $h_B(t) = \dfrac{2}{10}\left(\dfrac{t}{10}\right) = \dfrac{t}{50}$: at $t=5$, $h_B(5)=0.1$ (they match); at $t=15$, $h_B(15)=0.3 > h_A(15)=0.1$. In the short term both batches are comparable, but in the long term Batch B is less reliable.
		\item Batch A ($\beta=1$) does not age: it fails in a purely random (memoryless) fashion. Batch B ($\beta=2$) does age: its failure risk grows linearly with time, typical of accumulated mechanical wear. \quad \qedhere
	\end{enumerate}
\end{solucion}

\begin{solucion}
	\label{sol:4.7.9}
	For Problem \ref{prob:4.7.9}:
	\begin{enumerate}
		\item Substituting $\alpha=\nu/2$, $\beta=2$ in $M_X(t)=(1-\beta t)^{-\alpha}$ gives $M_X(t) = (1-2t)^{-\nu/2}$, which is exactly the MGF of $\chi^{2}_{\nu}$.
		\item Each $Z_i^{2} \sim \text{Gamma}(1/2,2) = \chi^{2}_{1}$. By the property of the sum of independent gamma variables with the same scale $\beta=2$, $\sum_{i=1}^{\nu} Z_i^{2} \sim \text{Gamma}\!\left(\nu\cdot\tfrac{1}{2},\,2\right) = \text{Gamma}(\nu/2,2) = \chi^{2}_{\nu}$.
		\item With $\nu=3$, $S=Z_1^{2}+Z_2^{2}+Z_3^{2}\sim\chi^{2}_{3}$. Since $7.815$ is by definition the $0.95$ quantile of $\chi^{2}_{3}$, it follows that $P(S>7.815) = 1-0.95 = 0.05$. \quad \qedhere
	\end{enumerate}
\end{solucion}

\begin{solucion}
	\label{sol:4.7.10}
	For Problem \ref{prob:4.7.10} with $T_i \sim \text{Exp}(0.5) = \text{Gamma}(1,2)$ i.i.d.:
	\begin{enumerate}
		\item By the additive property of the gamma distribution (same scale $\beta=2$), $T=\sum_{i=1}^{5}T_i \sim \text{Gamma}(5,2)$, that is, $T$ follows an Erlang distribution with $5$ stages and rate $0.5$.
		\item $\E(T) = \alpha\beta = 5\cdot 2 = 10$ years. $\Var(T) = \alpha\beta^{2} = 5\cdot 4 = 20$.
		\item $\text{Weibull}(\beta=1,\eta=2)$ coincides with $\text{Exp}(1/\eta) = \text{Exp}(0.5)$, which is precisely the distribution of each individual time $T_i$ between replacements, not that of the cumulative time $T$ until the fifth replacement. The two models coincide only when $\beta_{\text{Weibull}}=1$ because in that case the Weibull reduces to the exponential, the same family that generates each $T_i$; however, the cumulative time $T$ is a sum of exponentials (Gamma/Erlang), not a Weibull. \quad \qedhere
	\end{enumerate}
\end{solucion}
